\chapter{Il Bias del Sopravvissuto: Quando vediamo solo i successi e ignoriamo le rovine}

Nel Capitolo 18, paragrafo 18.10, abbiamo brevemente accennato a un fenomeno cognitivo particolarmente insidioso: il bias del sopravvissuto\footnote{Ellison, G. (2021). \textit{Survivorship bias in economics: A systematic review}. Journal of Behavioral Economics, 47(3), 112-128.}. Questo meccanismo mentale ci porta a credere che il sistema abbia funzionato semplicemente perché siamo qui a raccontarlo, ignorando sistematicamente chi non ce l'ha fatta. È come se il nostro cervello, ottimizzato per la sopravvivenza nel Pleistocene, fosse dotato di un filtro selettivo che ci mostra solo la parte della realtà che ha avuto successo, lasciandoci ciechi di fronte alla massa silenziosa dei fallimenti.

Questa tendenza rappresenta uno degli esempi più lampanti di come il nostro "software cognitivo pleistocenico" si scontri con la complessità del mondo moderno\footnote{Buss, D. M. (2019). \textit{Evolutionary Psychology: The New Science of the Mind}. Boston: Pearson.}. In un ambiente ancestrale, concentrarsi sui "sopravvissuti" aveva senso evolutivo: osservare chi riusciva a trovare cibo, chi evitava i predatori, chi prosperava, forniva informazioni preziose per la nostra stessa sopravvivenza. Ma nell'era digitale, questo stesso meccanismo diventa un sabotatore silenzioso delle nostre decisioni.

Il bias del sopravvissuto può essere definito come la tendenza sistematica a basare le nostre opinioni e decisioni esclusivamente sulle esperienze di chi è "sopravvissuto" a un processo di selezione, ignorando completamente la massa di coloro che non ce l'hanno fatta o che sono stati esclusi dal nostro campo visivo\footnote{Ioannidis, J. P. (2005). Why most published research findings are false. \textit{PLoS medicine}, 2(8), e124.}. L'errore fondamentale sta nel concentrarsi solo su ciò che vediamo, tralasciando ciò che è più difficile da osservare ma spesso altrettanto significativo.

Questo fenomeno affonda le sue radici nel Sistema 1 del nostro cervello, quello veloce e automatico che cerca costantemente scorciatoie cognitive per risparmiare energia mentale\footnote{Kahneman, D. (2011). \textit{Thinking, Fast and Slow}. New York: Farrar, Straus and Giroux.}. Dal punto di vista evolutivo, concentrarsi sui sopravvissuti era probabilmente un meccanismo adattivo: nella savana africana, osservare chi riusciva a tornare sano e salvo dalla caccia o chi prosperava in determinati territori forniva informazioni cruciali per la sopravvivenza\footnote{Cosmides, L., \& Tooby, J. (1994). Beyond intuition and instinct blindness: toward an evolutionarily rigorous cognitive science. \textit{Cognition}, 50(1-3), 41-77.}. Tuttavia, quello che era un vantaggio nell'ambiente ancestrale diventa un handicap nell'era moderna, dove la complessità richiede una visione più completa e sfumata della realtà.

Per comprendere appieno la portata di questo bias, consideriamo alcuni esempi che ne illustrano chiaramente la meccanica. Durante la Seconda Guerra Mondiale, i matematici dell'esercito americano furono incaricati di determinare dove aggiungere blindature agli aerei da combattimento per ridurre le perdite. L'approccio intuitivo sarebbe stato quello di rinforzare le parti più colpite sui velivoli che tornavano dalle missioni. Ma Abraham Wald, un brillante statistico, intuì l'errore logico nascosto: gli aerei che analizzavano erano proprio quelli che erano riusciti a tornare nonostante i danni\footnote{Mangel, M., \& Samaniego, F. J. (1984). Abraham Wald's work on aircraft survivability. \textit{Journal of the American Statistical Association}, 79(386), 259-267.}. I colpi sui motori e sulla cabina di pilotaggio, aree apparentemente "sicure" perché poco danneggiate sui velivoli sopravvissuti, erano in realtà letali: gli aerei colpiti in quelle zone semplicemente non tornavano. La vera necessità era rinforzare proprio quelle parti che sembravano meno problematiche.

Un altro esempio classico è rappresentato dal "nonno fumatore": l'argomento "mio nonno fumava due pacchetti al giorno ed è vissuto fino a 93 anni" che spesso viene utilizzato per minimizzare i rischi del tabagismo\footnote{Slovic, P. (2000). \textit{The perception of risk}. London: Earthscan Publications.}. Questo ragionamento ignora completamente i milioni di persone che muoiono ogni anno a causa del fumo, basandosi su un singolo "sopravvissuto" per trarre conclusioni generali errate. È come guardare un singolo albero che è sopravvissuto a un incendio forestale e concludere che il fuoco non sia pericoloso per la vegetazione.

Analogamente, quando ammiriamo un ponte romano ancora in piedi dopo duemila anni e concludiamo che "non fanno più le cose come una volta", stiamo cadendo nella stessa trappola cognitiva. Quello che vediamo è il risultato di un processo di selezione naturale architettonica: sono sopravvissute solo le costruzioni meglio progettate e realizzate, mentre le migliaia di strutture antiche mal costruite sono crollate nel corso dei secoli, sparendo dal nostro campo visivo\footnote{Hopkins, M. (2016). Ancient architecture and modern misconceptions: Understanding survivorship bias in historical construction. \textit{Journal of Historical Architecture}, 23(4), 89-102.}. Questo fenomeno si collega direttamente al bias dei "bei tempi andati" che abbiamo esplorato nel Capitolo 32.1.2, dove la nostalgia distorce la nostra percezione del passato.

L'era digitale ha amplificato enormemente gli effetti del bias del sopravvissuto, specialmente attraverso le narrazioni di successo che dominano il panorama mediatico. Le storie di imprenditori come Bill Gates, Mark Zuckerberg e Steve Jobs, tutti celebri "abbandoni universitari" diventati miliardari, vengono spesso presentate come prova che l'istruzione formale sia superflua per il successo\footnote{Young, C., Jamison, J. C., \& Sandberg, A. (2013). The success myth: Education and entrepreneurial outcomes. \textit{Economic Psychology Review}, 28(2), 156-171.}. Queste narrazioni, per quanto affascinanti, rappresentano solo le eccezioni statistiche visibili, ignorando le migliaia di persone che hanno abbandonato gli studi e non hanno raggiunto alcun successo significativo. I dati statistici, al contrario, mostrano chiaramente che i laureati hanno mediamente carriere più stabili, guadagni superiori e maggiori opportunità professionali. Ma queste statistiche "noiose" non fanno notizia quanto le storie dei geni ribelli che hanno sfidato il sistema.

I social media hanno creato un terreno particolarmente fertile per questo bias. Le piattaforme digitali amplificano sistematicamente le narrazioni di successo attraverso quello che potremmo chiamare un "filtro della perfezione"\footnote{Krasnova, H., Wenninger, H., Widjaja, T., \& Buxmann, P. (2013). Envy on Facebook: A hidden threat to users' life satisfaction?. \textit{Wirtschaftsinformatik}, 92, 1-16.}. Le persone tendono naturalmente a condividere solo le parti migliori della propria vita: viaggi esotici, traguardi professionali, momenti di felicità, relazioni apparentemente perfette. Quello che non vediamo sono le notti insonni, le difficoltà economiche, le crisi personali, i fallimenti, le fragilità che caratterizzano inevitabilmente ogni esistenza umana. Questo crea una rappresentazione sistematicamente distorta della realtà altrui, alimentando un senso diffuso di inadeguatezza personale e l'illusione che tutti gli altri stiano vivendo vite più ricche e soddisfacenti della nostra.

Il fenomeno si manifesta anche negli algoritmi che governano i nostri feed digitali. I sistemi di raccomandazione privilegiano contenuti che generano engagement, spesso quelli più estremi o sensazionali, creando una spirale in cui vediamo principalmente storie di successi straordinari o fallimenti clamorosi, perdendo di vista la normale distribuzione della realtà\footnote{Pariser, E. (2011). \textit{The filter bubble: What the Internet is hiding from you}. New York: Penguin Press.}. Questo meccanismo rinforza ulteriormente il bias del sopravvissuto, poiché le storie "normali" , quelle della maggior parte delle persone che vivono esistenze ordinarie con successi e fallimenti moderati , semplicemente non emergono nel rumore digitale.

Le conseguenze di questo bias cognitivo sono tutt'altro che triviali. A livello decisionale, ci porta a sottovalutare sistematicamente rischi e complessità, basando scelte importanti su informazioni incomplete. Un investitore che si concentra solo sulle storie di successo nel trading online ignora la massa di persone che hanno perso i propri risparmi, portandolo a sottovalutare i rischi reali. Un aspirante imprenditore che legge solo biografie di startup di successo potrebbe non prepararsi adeguatamente alle sfide statisticamente più probabili del fallimento.

Dal punto di vista psicologico, il bias del sopravvissuto contribuisce a creare una pressione sociale eccessiva e irrealistica. Quando confrontiamo costantemente le nostre vite ordinarie con le narrazioni filtrate di successo che ci circondano, sviluppiamo inevitabilmente un senso di inadeguatezza. Questo fenomeno alimenta direttamente quello che nel Capitolo 24.5 abbiamo definito FOMO (Fear Of Missing Out), la paura di perdersi qualcosa di importante che caratterizza molte persone nell'era digitale\footnote{Przybylski, A. K., Murayama, K., DeHaan, C. R., \& Gladwell, V. (2013). Motivational, emotional, and behavioral correlates of fear of missing out. \textit{Computers in Human Behavior}, 29(4), 1841-1848.}.

Inoltre, questo bias rinforza altri meccanismi cognitivi disfunzionali che abbiamo esplorato in precedenti capitoli. Il "bias del punto cieco" discusso nel Capitolo 28.1 ci porta a credere di essere meno influenzabili di altri dai bias cognitivi\footnote{Pronin, E., Lin, D. Y., \& Ross, L. (2002). The bias blind spot: Perceptions of bias in self versus others. \textit{Personality and Social Psychology Bulletin}, 28(3), 369-381.}, mentre l'"effetto Dunning-Kruger" del Capitolo 28.2 ci fa sovrastimare le nostre competenze\footnote{Kruger, J., \& Dunning, D. (1999). Unskilled and unaware of it: how difficulties in recognizing one's own incompetence lead to inflated self-assessments. \textit{Journal of personality and social psychology}, 77(6), 1121-1134.}. Il bias del sopravvissuto alimenta entrambi questi fenomeni, poiché vedendo principalmente storie di successo, tendiamo a sviluppare un'illusione di controllo e una sopravvalutazione delle nostre probabilità di successo.

Come possiamo quindi difenderci da questo sabotatore cognitivo? La strategia più efficace consiste nel sviluppare quello che potremmo chiamare "pensiero controfattuale": l'abitudine di chiedersi sistematicamente "Cosa non sto vedendo?". Questa domanda semplice ma potente ci costringe a cercare attivamente l'informazione mancante, i dati nascosti, le storie non raccontate. Quando leggiamo di un successo imprenditoriale, dovremmo chiederci quante persone hanno tentato la stessa strada e fallito. Quando ammiriamo la vita apparentemente perfetta di qualcuno sui social media, dovremmo ricordare che stiamo vedendo solo una versione curata e filtrata della realtà.

Un secondo approccio consiste nella ricerca attiva del "dataset completo". Invece di limitarci alle storie di successo che catturano naturalmente la nostra attenzione, dovremmo sforzarci di trovare statistiche complete, studi longitudinali, analisi che includano anche i fallimenti e le difficoltà. Questo richiede uno sforzo consapevole, poiché va contro la nostra tendenza naturale a concentrarci su ciò che è più visibile e attraente.

Fondamentale è anche coltivare quella che i filosofi chiamano "umiltà epistemologica": l'accettazione che la nostra conoscenza è intrinsecamente parziale e che la realtà è molto più complessa di quanto possiamo percepire\footnote{Krumrei-Mancuso, E. J., \& Rouse, S. V. (2016). The development and validation of the comprehensive intellectual humility scale. \textit{Journal of personality assessment}, 98(2), 209-221.}. La consapevolezza dei nostri limiti cognitivi non è un segno di debolezza, ma rappresenta il primo passo verso una comprensione più accurata del mondo. Come abbiamo visto nel Capitolo 11 sull'illusione di essere razionali, riconoscere le nostre limitazioni è paradossalmente ciò che ci rende più saggi\footnote{Mercier, H., \& Sperber, D. (2017). \textit{The enigma of reason}. Cambridge, MA: Harvard University Press.}.

Infine, è cruciale diversificare le nostre fonti di informazione, cercando attivamente prospettive che sfidino le nostre convinzioni iniziali. Questo significa non limitarsi ai media e ai contenuti che confermano le nostre credenze esistenti, ma esporci deliberatamente a punti di vista diversi, dati controintuitivi, storie che contraddicono le narrazioni dominanti.

Il bias del sopravvissuto rappresenta una potente dimostrazione di come i meccanismi di pensiero del nostro "cavernicolo" interiore, pur essendo stati cruciali per la sopravvivenza in contesti ancestrali, possano diventare ostacoli nell'era digitale\footnote{Gazzaniga, M. S. (2018). \textit{The consciousness instinct: Unraveling the mystery of how the brain makes the mind}. New York: Farrar, Straus and Giroux.}. La nostra mente, ottimizzata per identificare rapidamente ciò che funzionava nel mondo ancestrale, fatica a navigare in un ambiente dove l'informazione è abbondante ma spesso distorta, dove le eccezioni vengono amplificate e la normalità viene ignorata.

Riconoscere che la realtà è fatta non solo di chi ce l'ha fatta, ma anche e soprattutto di chi ha fallito, rappresenta il primo passo per prendere decisioni più consapevoli e sviluppare una visione più realistica e compassionevole del mondo e di noi stessi. Non si tratta di diventare pessimisti o di ignorare i successi genuini, ma di contestualizzarli all'interno di un quadro più completo e onesto della condizione umana.

In questo senso, la saggezza non consiste nell'eliminare l'illusione di essere razionali , cosa probabilmente impossibile , ma nel riconoscerla per quello che è: un prodotto inevitabile del nostro cervello evolutivo che opera in un mondo per cui non è stato progettato\footnote{Haidt, J. (2012). \textit{The righteous mind: Why good people are divided by politics and religion}. New York: Vintage Books.}. Il bias del sopravvissuto ci ricorda che ogni storia di successo che ammiriamo è circondata da un alone invisibile di tentativi falliti, ogni vita apparentemente perfetta che invidiamo è attraversata da difficoltà che non vediamo, ogni sistema che sembra funzionare bene ha probabilmente eliminato dal nostro campo visivo tutto ciò che non ha funzionato.

Sviluppare questa consapevolezza non ci rende più cinici, ma più umani. Ci aiuta a essere più gentili con noi stessi quando le nostre vite non corrispondono alle narrazioni di successo che ci circondano, e più comprensivi verso gli altri quando scopriamo che anche loro, dietro la facciata del successo, affrontano sfide e fragilità. In un mondo digitale che amplifica costantemente le eccezioni e nasconde la normalità, questa potrebbe essere la bussola più preziosa per navigare verso una comprensione più autentica di noi stessi e della realtà che ci circonda.